\chapter{Input--output theory}
\label{ch:input-output}

The master equation of \cref{ch:open-systems} discards the bath
entirely, keeping only its effect on the system through a few rates.
That is exactly what we want when the environment is a featureless
sink. But often the ``bath'' is a transmission line that we
actively drive and measure: a microwave cable carrying a readout
tone into a cavity and the reflected or transmitted signal back out.
We then need to keep the propagating field, not trace it away ---
we need to relate the field \emph{going in} to the field
\emph{coming out}, with the system in between. Input--output theory
provides exactly this relation. We model the bath as a continuum of
travelling modes, derive the quantum Langevin equation for the
system operator, and obtain the boundary condition
$\hat a_{\rm out}=\hat a_{\rm in}+\sqrt{\kappa}\,\hat a$ that ties
the outgoing field to the incoming field and the intracavity
amplitude. From it we read off cavity reflection and transmission
coefficients --- the $S$-parameters measured in every cQED
experiment --- and the formalism that underlies dispersive readout
(Part~III) and the single-photon-counter efficiency of Part~VI.

% --------------------------------------------------------------
\section{The bath as a continuum of modes}
\label{sec:continuum-bath}

\begin{phenomenon}
A semi-infinite transmission line attached to a cavity is not a
structureless reservoir: it carries fields that propagate toward
the cavity (the \emph{input}) and away from it (the \emph{output}).
Treating the line as a continuum of harmonic modes, with the cavity
coupled to all of them over a band around its resonance, lets us
follow those propagating fields explicitly. The same coupling that,
after tracing out the line, produced the decay term $\kappa\,
\mathcal D[\hat a]$ in the master equation now appears as a
\emph{boundary condition} relating the fields at the cavity port.
\end{phenomenon}

\paragraph{The model Hamiltonian.}
Let the system have annihilation operator $\hat a$ and Hamiltonian
$\op H_{\rm sys}$. Couple it to a continuum of bath modes
$\hat b(\omega)$ with $\comm{\hat b(\omega)}{\hat b^\dagger(\omega')}=\delta(\omega-\omega')$,
\begin{equation}
\label{eq:io-hamiltonian}
\op H \;=\; \op H_{\rm sys}
\;+\; \hbar\!\int_0^\infty\!\!\dd\omega\;\omega\,\hat b^\dagger(\omega)\hat b(\omega)
\;+\; i\hbar\!\int_0^\infty\!\!\dd\omega\;\sqrt{\frac{\kappa}{2\pi}}
\Bigl[\hat b^\dagger(\omega)\hat a - \hat a^\dagger\hat b(\omega)\Bigr],
\end{equation}
where the \emph{first Markov approximation} has been used to take
the coupling $\sqrt{\kappa/2\pi}$ independent of frequency over the
relevant band (a flat, white reservoir). The constant $\kappa$ will
turn out to be the cavity energy-decay rate.

\paragraph{Input and output fields.}
Solving the Heisenberg equation for $\hat b(\omega)$ formally and
substituting back defines two fields built from the bath operators
evaluated at the initial time $t_0<t$ and the final time $t_1>t$:
\begin{equation}
\label{eq:io-fields}
\hat a_{\rm in}(t) = \frac{-1}{\sqrt{2\pi}}\!\int\!\dd\omega\,e^{-i\omega(t-t_0)}\hat b_0(\omega),
\qquad
\hat a_{\rm out}(t) = \frac{1}{\sqrt{2\pi}}\!\int\!\dd\omega\,e^{-i\omega(t-t_1)}\hat b_1(\omega).
\end{equation}
The input field is the freely propagating field incident on the
system (set by the source before it interacts); the output field is
the freely propagating field that leaves (what a detector
downstream measures). Both obey the free-field commutator
$\comm{\hat a_{\rm in}(t)}{\hat a_{\rm in}^\dagger(t')}=\delta(t-t')$,
i.e.\ they are white-noise quantum fields.

% --------------------------------------------------------------
\section{The quantum Langevin equation and the boundary relation}
\label{sec:langevin}

\begin{statement}[\textbf{Quantum Langevin equation and input--output relation}~\cite{gardiner1985,gardiner2004,clerk2010}]
\label{stmt:input-output}
Eliminating the bath from \cref{eq:io-hamiltonian} yields a
Heisenberg--Langevin equation for any system operator $\hat a$,
\begin{equation}
\label{eq:qle}
\dot{\hat a} \;=\; -\frac{i}{\hbar}\comm{\hat a}{\op H_{\rm sys}}
\;-\; \frac{\kappa}{2}\,\hat a \;+\; \sqrt{\kappa}\,\hat a_{\rm in}(t),
\end{equation}
in which the bath enters only through the damping $-\tfrac{\kappa}{2}\hat a$
and the driving quantum noise $\sqrt{\kappa}\,\hat a_{\rm in}$. The
input, system, and output fields at the port are tied by the
boundary relation
\begin{equation}
\label{eq:input-output-relation}
\boxed{\;\hat a_{\rm out}(t) \;=\; \hat a_{\rm in}(t) + \sqrt{\kappa}\,\hat a(t)\;}
\end{equation}
expressing energy conservation at the interface: what comes out is
what went in, plus what the system radiates.
\end{statement}

\paragraph{Consistency with the master equation.}
Taking the expectation of \cref{eq:qle} for a vacuum input
($\langle\hat a_{\rm in}\rangle=0$) gives
$\langle\dot{\hat a}\rangle=-\tfrac{i}{\hbar}\langle\comm{\hat a}{\op H_{\rm sys}}\rangle-\tfrac{\kappa}{2}\langle\hat a\rangle$,
which is exactly the equation of motion produced by the Lindblad
dissipator $\kappa\,\mathcal D[\hat a]$ of \cref{eq:dissipator}
(use $\langle\dot{\hat a}\rangle=\Tr[\hat a\,\dot{\op\rho}]$ with
the adjoint dissipator). Input--output theory and the master
equation describe the same open dynamics: the master equation
follows from tracing out $\hat a_{\rm in}$, while input--output
keeps it so we can drive and measure through the port.

\paragraph{Why the half.}
The amplitude decays at $\kappa/2$ in \cref{eq:qle} while the energy
$\langle\hat a^\dagger\hat a\rangle$ decays at $\kappa$ --- the same
factor-of-two between amplitude and population that appeared for the
qubit coherence in \cref{ch:relaxation-dephasing}. It is the generic
relation between a field amplitude decay rate and the corresponding
intensity (photon-number) decay rate.

% --------------------------------------------------------------
\section{Cavity response: reflection and transmission}
\label{sec:cavity-response}

\begin{phenomenon}
Drive a cavity through a port and the steady-state field it builds
up, and the field it returns, depend on how close the drive is to
resonance and on how strongly the cavity is coupled to the line.
The reflected (or transmitted) amplitude as a function of drive
frequency is the cavity's $S$-parameter --- the directly measured
quantity in a microwave reflection or transmission experiment, and
the carrier of the qubit-state information in dispersive readout.
\end{phenomenon}

\begin{statement}[\textbf{Single-port reflection coefficient}]
\label{stmt:reflection}
For an empty linear cavity, $\op H_{\rm sys}=\hbar\omega_c\hat a^\dagger\hat a$,
driven by a coherent input at frequency $\omega_d$, the steady-state
of \cref{eq:qle} gives the intracavity amplitude
\begin{equation}
\label{eq:intracavity}
\langle\hat a\rangle \;=\; \frac{\sqrt{\kappa}\,\langle\hat a_{\rm in}\rangle}{i(\omega_c-\omega_d) + \kappa/2},
\end{equation}
and, through the boundary relation \cref{eq:input-output-relation},
the reflection coefficient
\begin{equation}
\label{eq:reflection}
r(\omega_d) \;=\; \frac{\langle\hat a_{\rm out}\rangle}{\langle\hat a_{\rm in}\rangle}
\;=\; \frac{i(\omega_c-\omega_d) - \kappa/2}{i(\omega_c-\omega_d) + \kappa/2}.
\end{equation}
For a single lossless port $|r|=1$ at all frequencies: the cavity
returns all the power, but imprints a frequency-dependent
\emph{phase} that swings by $2\pi$ across the resonance, with the
steepest slope at $\omega_d=\omega_c$.
\end{statement}

\paragraph{Two ports: transmission and linewidth.}
With an input port (rate $\kappa_1$) and an output port (rate
$\kappa_2$), the total linewidth is $\kappa=\kappa_1+\kappa_2+\kappa_{\rm int}$
(including internal loss $\kappa_{\rm int}$), and the transmission
\begin{equation}
\label{eq:transmission}
t(\omega_d) \;=\; \frac{\sqrt{\kappa_1\kappa_2}}{i(\omega_c-\omega_d)+\kappa/2}
\end{equation}
is a Lorentzian of full width $\kappa$ centred at $\omega_c$, peak
height $2\sqrt{\kappa_1\kappa_2}/\kappa$. This is the $S_{21}$ trace
of resonator spectroscopy; fitting its centre and width yields
$\omega_c$ and the loaded quality factor $Q=\omega_c/\kappa$.

\paragraph{Side-coupled (notch) geometry.}
The readout resonators of circuit QED are usually placed not
\emph{inline} but \emph{beside} a continuous feedline: the line runs
straight from port~1 to port~2, and the cavity hangs off it through a
coupling capacitor (a \emph{hanger} or \emph{notch} resonator). The
feedline is now the propagating bath of \cref{eq:io-hamiltonian}; the
cavity radiates back into it at the external rate $\kappa_c$ and
dissipates internally at $\kappa_i$, so the loaded linewidth is
$\kappa=\kappa_c+\kappa_i$. The field transmitted \emph{past} the
coupling point is the through-going wave minus what the cavity
re-radiates forward, and the boundary relation
\cref{eq:input-output-relation} gives
\begin{equation}
\label{eq:notch-s21}
S_{21}(\omega_d) \;=\; 1 \;-\; \frac{\kappa_c/2}{\,i(\omega_c-\omega_d)+\kappa/2\,}
\;=\; 1 \;-\; \frac{Q_l/Q_c}{\,1+2iQ_l\,\delta/\omega_c\,},
\end{equation}
with $\delta=\omega_d-\omega_c$ and the quality factors
$Q_c=\omega_c/\kappa_c$ (external/coupling), $Q_i=\omega_c/\kappa_i$
(internal), and loaded $Q_l^{-1}=Q_c^{-1}+Q_i^{-1}$. Unlike the inline
Lorentzian \emph{peak} of \cref{eq:transmission}, the notch geometry
produces a transmission \emph{dip}: off resonance $S_{21}\to1$ (the
wave ignores the cavity), while on resonance the cavity scatters and
$S_{21}(\omega_c)=1-Q_l/Q_c=Q_i/(Q_i+Q_c)$, which vanishes for a
lossless, strongly overcoupled resonator ($\kappa_c\gg\kappa_i$).
Fitting the complex dip in the $(\Re,\Im)$ plane yields $\omega_c$,
$Q_l$, and the coupling/internal split~\cite{gao2008,probst2015}. We
will meet \emph{exactly} this functional form again in
\cref{sec:waveguide-readout}, where the scatterer is not a linear
cavity but a single qubit coupled straight to the line.

\paragraph{Application: dispersive readout.}
In the dispersive regime (Part~III), the qubit shifts the cavity
frequency by $\pm\chi$ depending on its state, $\omega_c\to\omega_c\pm\chi$.
Through \cref{eq:reflection}, the two qubit states produce two
different reflected phases; measuring that phase with a
near-resonant tone is a quantum non-demolition readout of the qubit.
The phase contrast is largest when $2\chi\sim\kappa$, the matching
condition behind the photon-resolution requirement quoted in
Part~VI. The same input--output bookkeeping fixes the detection
efficiency $\eta$ of the single-microwave-photon counter
(\cref{ch:microwave-absorption}): the matched condition there,
$\kappa_b=\kappa_{\rm nl}$, is precisely the impedance match that
maximises power transfer through the port.

\paragraph{Transition.}
The boundary relation \cref{eq:input-output-relation} treats one
system and one port. Chaining the output of one system into the
input of the next describes a measurement chain --- a cavity feeding
an amplifier feeding a detector --- and provides the launch point
for the continuous-measurement description of the next chapter.

% --------------------------------------------------------------
\section{Cascaded systems and the route to continuous measurement}
\label{sec:cascaded}

\begin{statement}[\textbf{Cascaded input--output}]
\label{stmt:cascaded}
When the output of system $1$ is fed unidirectionally into the
input of system $2$ (e.g.\ through a circulator), the input of the
second is the output of the first, delayed by the propagation time
$\tau$:
\begin{equation}
\label{eq:cascade}
\hat a_{\rm in}^{(2)}(t) \;=\; \hat a_{\rm out}^{(1)}(t-\tau).
\end{equation}
In the coarse-grained (Markovian) limit $\tau\to0$, this chains the
Langevin equations into a single cascaded master equation for the
joint system, with a one-way (chiral) coupling that has no
back-action from $2$ onto $1$.
\end{statement}

\paragraph{From the output field to a measurement record.}
A homodyne or photon-counting detector placed in the output field
$\hat a_{\rm out}$ converts the propagating quantum field into a
classical record. Because $\hat a_{\rm out}$ carries, through
\cref{eq:input-output-relation}, information about the system
operator $\hat a$, monitoring it continuously back-acts on the
system state. Conditioning the system state on the measured record
turns the deterministic master equation into a \emph{stochastic}
one --- the quantum-trajectory description that
\cref{ch:continuous-measurement} develops. Input--output theory is
thus the bridge between the system and the measuring apparatus:
the master equation is what remains after averaging over all
possible output records.

% --------------------------------------------------------------
This chapter kept the bath that \cref{ch:open-systems} discarded,
modelling it as a continuum of propagating modes. The quantum
Langevin equation \cref{eq:qle} and the boundary relation
$\hat a_{\rm out}=\hat a_{\rm in}+\sqrt{\kappa}\,\hat a$ relate the
incoming and outgoing fields through the system, reproduce the
master-equation damping on average, and yield the cavity reflection
and transmission coefficients measured in the laboratory. The same
relations underlie dispersive readout and set the efficiency of the
photon counters in Part~VI.

The output field carries the system's information to the detector;
conditioning the state on a continuous measurement of that field is
the final topic of Part~II. \Cref{ch:continuous-measurement} takes
up the quantum Zeno effect and the unravelling of the master
equation into individual quantum trajectories.

\clearpage
\begin{chaptersummary}

\paragraph{Continuum bath.}
A transmission line is modelled as a continuum of travelling modes
coupled to the system with a flat rate $\sqrt{\kappa/2\pi}$ (first
Markov approximation). The freely propagating incident and emitted
fields are the input and output fields, white-noise operators with
$\comm{\hat a_{\rm in}(t)}{\hat a_{\rm in}^\dagger(t')}=\delta(t-t')$.

\paragraph{Quantum Langevin and input--output.}
\begin{equation}
\dot{\hat a} = -\frac{i}{\hbar}\comm{\hat a}{\op H_{\rm sys}}
- \frac{\kappa}{2}\hat a + \sqrt{\kappa}\,\hat a_{\rm in},
\qquad
\hat a_{\rm out} = \hat a_{\rm in} + \sqrt{\kappa}\,\hat a.
\end{equation}
On average this reproduces the master-equation damping
$\kappa\mathcal D[\hat a]$; the amplitude decays at $\kappa/2$, the
energy at $\kappa$.

\paragraph{Cavity response.}
Single-port reflection
$r=\dfrac{i(\omega_c-\omega_d)-\kappa/2}{i(\omega_c-\omega_d)+\kappa/2}$,
$|r|=1$ for a lossless port with a $2\pi$ phase swing across
resonance; two-port transmission is a Lorentzian of width
$\kappa=\kappa_1+\kappa_2+\kappa_{\rm int}$ ($S_{21}$, loaded
$Q=\omega_c/\kappa$). A state-dependent cavity pull $\pm\chi$ is the
basis of dispersive readout.

\paragraph{Cascade and measurement.}
Feeding one output into the next input (delayed) chains systems
into a cascaded, one-way master equation. Detecting the output
field continuously back-acts on the system and yields the quantum
trajectories of \cref{ch:continuous-measurement}; the master
equation is the average over all records.

\end{chaptersummary}

% --------------------------------------------------------------
\section*{Exercises}
\addcontentsline{toc}{section}{Exercises}

\begin{exercise}[Free-field commutator of the input]
Using $\comm{\hat b(\omega)}{\hat b^\dagger(\omega')}=\delta(\omega-\omega')$
and the definition \cref{eq:io-fields}, show that
$\comm{\hat a_{\rm in}(t)}{\hat a_{\rm in}^\dagger(t')}=\delta(t-t')$.
\begin{solution}
$\comm{\hat a_{\rm in}(t)}{\hat a_{\rm in}^\dagger(t')}
=\tfrac1{2\pi}\!\int\!\dd\omega\dd\omega'\,e^{-i\omega(t-t_0)}e^{i\omega'(t'-t_0)}
\comm{\hat b_0(\omega)}{\hat b_0^\dagger(\omega')}
=\tfrac1{2\pi}\!\int\!\dd\omega\,e^{-i\omega(t-t')}=\delta(t-t')$.
The input field is delta-correlated white noise --- the time-domain
signature of the flat (Markovian) reservoir.
\end{solution}
\end{exercise}

\begin{exercise}[Energy versus amplitude decay]
From the vacuum-input Langevin equation, show that
$\langle\hat a\rangle$ decays at $\kappa/2$ while
$\langle\hat a^\dagger\hat a\rangle$ decays at $\kappa$ for an empty
cavity. Reconcile with the master-equation result.
\begin{solution}
With $\op H_{\rm sys}=\hbar\omega_c\hat a^\dagger\hat a$ and vacuum
input, $\langle\dot{\hat a}\rangle=(-i\omega_c-\kappa/2)\langle\hat a\rangle$,
so $|\langle\hat a\rangle|\propto e^{-\kappa t/2}$. For the
population, $\tfrac{\dd}{\dd t}\langle\hat a^\dagger\hat a\rangle
=-\kappa\langle\hat a^\dagger\hat a\rangle$ from $\kappa\mathcal D[\hat a]$,
so the energy decays at $\kappa$. Amplitude rate is half the energy
rate, exactly as for qubit coherence vs population in
\cref{ch:relaxation-dephasing}.
\end{solution}
\end{exercise}

\begin{exercise}[Steady-state intracavity field]
Derive \cref{eq:intracavity} by setting $\langle\dot{\hat a}\rangle=0$
in the driven Langevin equation for a coherent input at $\omega_d$
(work in the frame rotating at $\omega_d$).
\begin{solution}
In the rotating frame $\op H_{\rm sys}\to\hbar(\omega_c-\omega_d)\hat a^\dagger\hat a$,
so $\langle\dot{\hat a}\rangle=[-i(\omega_c-\omega_d)-\kappa/2]\langle\hat a\rangle+\sqrt\kappa\langle\hat a_{\rm in}\rangle$.
Setting this to zero gives
$\langle\hat a\rangle=\sqrt\kappa\langle\hat a_{\rm in}\rangle/[i(\omega_c-\omega_d)+\kappa/2]$,
which is \cref{eq:intracavity}.
\end{solution}
\end{exercise}

\begin{exercise}[Reflection phase and its slope]
For the single-port reflection \cref{eq:reflection}, confirm $|r|=1$
and compute the phase $\arg r$ as a function of detuning. Find the
slope $\dd(\arg r)/\dd\omega_d$ at resonance.
\begin{solution}
Numerator and denominator are complex conjugates up to sign:
$|i\delta-\kappa/2|=|i\delta+\kappa/2|=\sqrt{\delta^2+\kappa^2/4}$
with $\delta=\omega_c-\omega_d$, so $|r|=1$. The phase is
$\arg r=\pi-2\arctan(2\delta/\kappa)$ (mod $2\pi$), swinging by
$2\pi$ across resonance. At $\delta=0$,
$\dd(\arg r)/\dd\delta=-4/\kappa$, so
$\dd(\arg r)/\dd\omega_d=+4/\kappa$ --- the steepest phase response,
hence the best frequency discrimination, at resonance and for small
$\kappa$.
\end{solution}
\end{exercise}

\begin{exercise}[Lorentzian transmission and $Q$]
From \cref{eq:transmission} show that $|t(\omega_d)|^2$ is a
Lorentzian of full width at half maximum $\kappa$, and relate the
loaded quality factor to $\kappa$.
\begin{solution}
$|t|^2=\kappa_1\kappa_2/[(\omega_c-\omega_d)^2+\kappa^2/4]$, a
Lorentzian peaked at $\omega_d=\omega_c$ with half-maximum where
$(\omega_c-\omega_d)^2=\kappa^2/4$, i.e.\ FWHM $=\kappa$. The loaded
quality factor is $Q=\omega_c/\kappa$, the ratio of resonance
frequency to linewidth.
\end{solution}
\end{exercise}

\begin{exercise}[Critical coupling]
For a cavity with one external port (rate $\kappa_1$) and internal
loss $\kappa_{\rm int}$, the reflection is
$r=\dfrac{i\delta-(\kappa_1-\kappa_{\rm int})/2}{i\delta+(\kappa_1+\kappa_{\rm int})/2}$.
Show that on resonance the reflected power vanishes when
$\kappa_1=\kappa_{\rm int}$ (critical coupling) and interpret.
\begin{solution}
At $\delta=0$, $r=-(\kappa_1-\kappa_{\rm int})/(\kappa_1+\kappa_{\rm int})$,
which is zero iff $\kappa_1=\kappa_{\rm int}$. At critical coupling
all incident power is absorbed by the internal loss --- the external
coupling is impedance-matched to the internal dissipation, leaving
no reflection. Over- ($\kappa_1>\kappa_{\rm int}$) and under-coupled
cases give non-zero reflection of opposite sign.
\end{solution}
\end{exercise}

\begin{exercise}[Dispersive cavity pull]
With a state-dependent cavity frequency $\omega_c\pm\chi$, compute
the difference in reflected phase at $\omega_d=\omega_c$ between the
two qubit states for $2\chi\ll\kappa$ and for $2\chi=\kappa$. Which
gives better state discrimination?
\begin{solution}
Using $\arg r=\pi-2\arctan(2\delta/\kappa)$ with $\delta=\pm\chi$:
the two phases are $\pi\mp2\arctan(2\chi/\kappa)$, separated by
$\Delta\phi=4\arctan(2\chi/\kappa)$. For $2\chi\ll\kappa$,
$\Delta\phi\approx8\chi/\kappa$ (small, poor contrast); for
$2\chi=\kappa$, $\Delta\phi=4\arctan1=\pi$ (maximal contrast). So
$2\chi\sim\kappa$ is the matched condition for readout, as quoted in
Part~VI.
\end{solution}
\end{exercise}

\begin{exercise}[Master equation from the Langevin equation]
Show that averaging \cref{eq:qle} over a vacuum input reproduces the
equation of motion $\langle\dot{\hat a}\rangle=\Tr[\hat a\,\dot{\op\rho}]$
with $\dot{\op\rho}$ given by the Lindblad dissipator
$\kappa\mathcal D[\hat a]$.
\begin{solution}
The adjoint dissipator acting on $\hat a$ is
$\kappa(\hat a^\dagger\hat a\hat a-\tfrac12\{\hat a^\dagger\hat a,\hat a\})
=\kappa(\hat a^\dagger\hat a\hat a-\tfrac12\hat a^\dagger\hat a\hat a-\tfrac12\hat a\hat a^\dagger\hat a)
=\tfrac{\kappa}2(\hat a^\dagger\hat a\hat a-\hat a\hat a^\dagger\hat a)
=\tfrac\kappa2\comm{\hat a^\dagger}{\hat a}\hat a=-\tfrac\kappa2\hat a$.
Adding the unitary part gives
$\langle\dot{\hat a}\rangle=-\tfrac i\hbar\langle\comm{\hat a}{\op H_{\rm sys}}\rangle-\tfrac\kappa2\langle\hat a\rangle$,
the vacuum-average of \cref{eq:qle}.
\end{solution}
\end{exercise}

\begin{exercise}[One-way cascade has no back-action]
Argue from \cref{eq:cascade} why, in a cascaded chain with a
circulator, system $2$ exerts no back-action on system $1$, and why
this requires the non-reciprocal element.
\begin{solution}
The relation $\hat a_{\rm in}^{(2)}=\hat a_{\rm out}^{(1)}$ feeds
$1$'s output into $2$ but provides no path for $2$'s output to reach
$1$'s input. The dynamics of $1$ depend only on its own input, so
$2$ cannot back-act. This unidirectionality is enforced physically
by the circulator/isolator (a non-reciprocal element); without it
the line is bidirectional and $2$'s emission would re-enter $1$.
\end{solution}
\end{exercise}

\begin{exercise}[Output field of an empty driven cavity]
Find $\langle\hat a_{\rm out}\rangle$ for the driven empty cavity in
steady state and verify $|\langle\hat a_{\rm out}\rangle|=|\langle\hat a_{\rm in}\rangle|$
for a lossless single port.
\begin{solution}
$\langle\hat a_{\rm out}\rangle=\langle\hat a_{\rm in}\rangle+\sqrt\kappa\langle\hat a\rangle
=\langle\hat a_{\rm in}\rangle[1+\kappa/(i\delta+\kappa/2)]
=\langle\hat a_{\rm in}\rangle\,(i\delta-\kappa/2)/(i\delta+\kappa/2)=r\langle\hat a_{\rm in}\rangle$
with $|r|=1$ (Exercise on reflection phase). Hence
$|\langle\hat a_{\rm out}\rangle|=|\langle\hat a_{\rm in}\rangle|$:
a lossless port returns all the power, consistent with energy
conservation in \cref{eq:input-output-relation}.
\end{solution}
\end{exercise}
